\documentclass{article}
\usepackage[backend=biber,natbib=true,style=authoryear]{biblatex}
\addbibresource{/home/nguyen/1_NQBH/math/bib.bib}
\usepackage[utf8]{inputenc}
\usepackage{float}
\usepackage{graphicx}
\usepackage[colorlinks=true,linkcolor=blue,urlcolor=red,citecolor=magenta]{hyperref}
\usepackage{amsmath,amssymb,amsthm}
\allowdisplaybreaks
\numberwithin{equation}{section}
\newtheorem{assumption}{Assumption}[section]
\newtheorem{lemma}{Lemma}[section]
\newtheorem{corollary}{Corollary}[section]
\newtheorem{definition}{Definition}[section]
\newtheorem{proposition}{Proposition}[section]
\newtheorem{theorem}{Theorem}[section]
\newtheorem{notation}{Notation}[section]
\newtheorem{remark}{Remark}[section]
\newtheorem{example}{Example}[section]
\newtheorem{ques}{Question}[section]
\newtheorem{problem}{Problem}[section]
\newtheorem{conjecture}{Conjecture}[section]
\usepackage[left=0.5in,right=0.5in,top=1.5cm,bottom=1.5cm]{geometry}
\def\labelitemii{$\circ$}

\title{don Miguel Ruiz, don Jose Ruiz (2011) (with Janet Mills). \textit{The Fifth Agreement: A Practical Guide to Self-Mastery} (A Toltec Wisdom Book)}
\author{Collector: Nguyen Quan Ba Hong}
\date{\today}

\begin{document}
\maketitle
\setcounter{secnumdepth}{6}
\setcounter{tocdepth}{6}
\tableofcontents

%------------------------------------------------------------------------------%

\begin{quotation}
    To every human who lives on this beautiful planet, and to the generations to come.
\end{quotation}

\section*{The Toltec}
\textsc{Thousands of years ago, the Toltec were} known throughout southern Mexico as ``women and men of knowledge.''

Anthropologists have spoken of the Toltec as a nation or a race, but, in fact, the Toltec were scientists and artists who formed a society to explore and conserve the spiritual knowledge and practices of the ancient ones.

They came together as masters (\textit{naguals}) and students at Teotihuacan, the ancient city of pyramids outside Mexico City known as the place where ``Man Becomes God.''

Over the millenia, the \textit{naguals} were forced to conceal the ancestral wisdom and maintain its existence in obscurity.

European conquest, coupled with rampant misuse of personal power by a few of the apprentices, made it necessary to shield the knowledge from those who were not prepared to use it wisely or who might intentionally misuse it for personal gain.

%
Fortunately, the esoteric Toltec knowledge was embodied and passed on through generations by different lineages of \textit{naguals}.

Though it remained veiled in secrecy for hundreds of years, ancient prophecies foretold the coming of an age when it would be necessary to return the wisdom to the people.

Now, don Miguel Ruiz and don Jose Ruiz (\textit{naguals} from the Eagle Knight lineage) have been guided to share with us the powerful teachings of the Toltec.

%
Toltec wisdom arises from the same essential unity of truth as all the sacred esoteric traditions found around the world.

Though it is not a religion, it honors all the spiritual masters who have taught on the earth.

While it does embrace spirit, it is most accurately described as \textit{a way of life}, distinguished by the \textit{ready accessibility of happiness and love}.

%------------------------------------------------------------------------------%

\section*{Introduction}
\begin{center}
    \large by don Miguel Ruiz
\end{center}
\textsc{The 4 Agreements was published many years} ago.

If you have read the book, you already know what these agreements can do.

They have the ability to transform your life by breaking thousands of limiting agreements you have made with yourself, with other people, with \textit{life} itself.

%
The 1st time you read \textit{The 4 Agreements}, it begins to work its magic.

It goes much deeper than the words you are reading.

You feel that you already know every word in the book.

You feel it, but perhaps you never put it into words.

When you read the book the 1st time, it challenges what you believe, and takes you to the limit of your comprehension.

You break many limiting agreements, and overcome many challenges, but then you see new challenges.

When you read the book a 2nd time, it feels as if you're reading a completely different book because the limits of your comprehension have already grown.

Once again, it takes you into a deeper awareness of yourself, and you reach the limit that you can reach in that moment.

And when you read the book a 3rd time, it's just as if you're reading another book.

%
Just like magic, because they \textit{are} magic, the 4 Agreements slowly help you to recover your authentic self.

With practice, these 4 simple agreements take you to what you \textit{really} are, not what you pretend to be, and this is exactly where you want to be: what you really are.

%
The principles in \textit{The 4 Agreements} speak to the heart of all human beings, from the young to the old.

They speak to people of different cultures all around the world - people who speak different languages, people whose religious and philosophical beliefs are vastly different.

They have been taught at different kinds of schools, from elementary schools to secondary schools to universities.

The principles in \textit{The 4 Agreements} reach everyone because they are \textit{pure common sense}.

%
Now it's time to give another gift: \textit{The 5th Agreement}.

The 5th agreement wasn't included in my 1st book because the 1st 4 agreements were enough of a challenge at that time.

The 5th agreement is made with words, of course, but its meaning and intent is beyond the words.

The 5th agreement is ultimately about seeing your whole reality with the eyes of truth, \textit{without} words.

The result of practicing the 5th agreement is the \textit{complete acceptance} of yourself just the way you are, and the complete acceptance of everybody else just the way they are.

\fbox{\textit{The reward is your eternal happiness}.}

%
Many years ago, I began teaching some of the concepts in this book to my apprentices, but then I stopped because nobody seemed to understand what I was trying to say.

Though I had shared the 5th agreement with my apprentices, I discovered that nobody was ready to learn the teachings that underlie this agreement.

Years later, my son, don Jose, started to share the same teachings with a group of students, and he succeeded where I had failed.

Maybe the reason don Jose was successful was because he had complete faith in sharing the message.

His very presence spoke the truth and challenged the beliefs of the people who attended his classes.

He made a huge difference in their lives.

%
Don Jose Ruiz has been my apprentice since he was a child, since he learned to speak.

In this book, I [don Miguel Ruiz] honored to introduce my son, and to present the essence of the teachings we delivered together over a period of 7 years.

%
To keep the message as personal as possible, and in keeping with the 1st-person voice of prior books in the Toltec Wisdom series, we have opted to present  \textit{The 5th Agreement} in the same 1st-person style of writing.

In this book, we speak to the reader with 1 voice, and with 1 heart.

%------------------------------------------------------------------------------%

\begin{center}
    \huge\textsc{Part I: The Power of Symbols}
\end{center}

\section{In the Beginning: It's All in the Program}
\textsc{From the moment you are born, you deliver} a message to the world.

\textit{What is the message?}

The message is \textit{you}, that child.

It's the presence of an \textit{angel}, a messenger from the infinite in a human body.

The infinite, a total power, creates a program just for you, and everything you need to be what you are is in the program.

You are born, you grown up, you mate, you grow old, and in the end you return to the infinite.

Every cell in your body is a universe of its own.

\textit{It's intelligent, it's complete, and it's programmed to be whatever it is}.

%
You are programmed to be \textit{you}, whatever you are, and it makes no difference to the program what you mind \textit{thinks} you are.

The program is not in the thinking mind.

It's in the body, in what we call the \textit{DNA}, and in the beginning, you instinctively follow its wisdom.

As a very young child, you know what you like, what you don't like, when you like it, when you don't.

You follow what you like, and you try to avoid what you don't like.

You follow your instincts, and those instinct guide you to be happy, to enjoy life, to play, to love, to fulfill your needs.

Then what happens?

Your body begins to develop, your mind begins to mature, and you begin to use symbols to deliver your message.

Just as the birds understand the birds, and the cats understands the cats, the humans understand the humans through a symbology.

If you were born on an island and then lived all alone, it might take you 10 years, but you would give a name to everything that you see, and you would use that language to communicate a message, even if it was only to yourself.

What would you do this?

Well, it's easy to understand, and it's not because humans are so intelligent.

It's because we are programmed to create a language, to invent an entire symbology for ourselves.

%
As you know, all around the world humans speak and write in thousands of different languages.

Humans have invented all kinds of symbols to communicate not only with other humans but more importantly with ourselves.

The symbols can be sounds that we speak, motions that we make, or handwriting and signs that are graphic in nature.

There are symbols for objects, ideas, music, and mathematics, but the introduction of sounds is the very 1st step, which means we learn to use symbols to speak.

%
The humans who come before us already have names for everything that exists, and they teach us the meaning of sounds.

They call this a \textit{table}; they call that a \textit{chair}.

They also have names for things that only exist in our imagination, like mermaids and unicorns.

Every word that we learn is a symbol for something real or imagined, and there are thousands of words to learn.

If we observe children who are 1--4 years old, we can see the effort they make trying to learn an entire symbology.

It's a big effort that we usually don't remember because our mind is not yet mature, but with repetition and practice we finally learn to speak.

%
Once we learn to speak, the humans who take care of us teach us what they know, which means they program us with knowledge.

The humans we live with have lots of knowledge, including all the social, religious, and moral rules of their culture.

They hook our attention, pass on the information, and teach us to be like them.

We learn how to be a man or woman according to the society in which we are born.

We learn how to behave the ``right'' way in our society, which means how to be a ``good'' human.

%
In truth, we are domesticated the same way that a dog, a cat, or any animal is domesticated: through a system of punishment and reward.

We are told that we're a \textit{good boy} or a \textit{good girl} when we do what the grown-ups want us to do; we're a \textit{bad boy} or a \textit{bad girl} when we don't do what they want us to do.

Sometimes we are punished without being bad, and sometimes we are rewarded without being good.

Out of fear of being punished and fear of not getting a reward, we start trying to please other people.

We try to be good, because bad people don't receive rewards; they are punished.

%
In human domestication, all the rules and values of our family and society are imposed on us.

We don't have the opportunity to choose our beliefs; we are told what to believe, and what not to believe.

The people we live with tell us their opinions: what is good and what is bad, what is right and what is wrong, what is beautiful and what is ugly.

Just like a computer, all that information is downloaded into our head.

We are innocent; we \textit{believe} what our parents or other grown-ups tell us; we \textit{agree}, and the information is stored in our memory.

Everything we learn goes into our mind by agreement, and it stays in our mind by agreement, but 1st it goes through the attention.

%
The attention is very important in humans because it's the part of the mind that makes it possible for us to concentrate on a single object or thought out of a whole range of possibilities.

\textit{Through the attention, information from the outside is conveyed to the inside and vice versa}.

The attention is the channel we use to send and receive messages from human to human.

It's like \textit{a bridge from 1 mind to another mind}; we open the bridge with sounds, signs, symbols, touch - with any event that hooks the attention.

This is how we teach, and this is how we learn.

We cannot teach anything if we don't have someone's attention; we cannot learn anything if we don't pay attention.

%
Using the attention, the grown-ups teach us how to create an entire reality in our mind with the use of symbols.

After they teach us a symbology by sound, the grown-ups drill us with our ABCs, and we learn the same language, but graphically.

Our imagination begins to develop, our curiosity grows stronger, and we start to ask questions.

We ask and ask, and we keep asking questions; we gather information from everywhere.

And we know that we've finally mastered a language when we are able to use the symbols to talk to ourselves in our head.

This is when we learn to \textit{think}.

Before that, we don't think; we mimic sounds and use symbols to communicate, but life is simple before we attach any meaning or emotion to the symbols.

%
Once we give meaning to the symbols, we being to use them to try to make sense of everything that happens in our lives.

We use the symbols to think about things that are real, and to think about things that aren't real, but that we start to imagine are real, like beautiful and ugly, skinny and fat, smart and stupid.

And if you notice, \textit{we can only think in a language that we master}.

For many years, I spoke only Spanish, and it took a long time for me to master enough symbols in English to think in English.

To master a language is not easy, but at a certain point, we find ourselves \textit{thinking} with the symbols we learn.

%
By the time we go to school, when we are 5 or 6 years old, we understand the meaning of abstract concepts like right and wrong, winter and loser, perfect and imperfect.

In school, we continue to learn how to read and write the symbols we already know, and the written language makes it possible for us to accumulate more knowledge.

We continue to give meaning to more and more symbols, and thinking becomes not only effortless but automatic.

%
Now the symbols that we learned are hooking our attention all by themselves.

It's what we know that's talking to us, and we are listening to what our knowledge says.

I call it \fbox{\textit{the voice of knowledge}} because \textit{knowledge is talking in our head}.

Many times we hear the voice with different tonalities; we hear the voice of our mother, our father, our brothers and sisters, and the voice never stops talking.

\textit{The voice isn't real; it's our creation}.

But we \textit{believe} that it's real because we give it life through the power of our faith, which means we believe \textit{without a doubt} what that voice is telling us.

This is when the opinions of the humans around us start taking over our mind.

%
Everybody has an opinion of us, and they tell us what we are.

As very young children, we don't know what we are.

The only way we can see ourselves is through a mirror, and other people act as that mirror.

Our mother tells us what we are, and we believe her.

It's completely different from what our father tells us, or what our brothers and sisters tell us, but we agree with them, too.

People tell us how we look, and it's especially true when we are little children.

``Look, you have the eyes of your mother, the nose of your grandfather.''

We listen to all the opinions of our family, our teachers, and the big children at school.

We see our image in those mirrors, we agree that this is what we are, and as soon as we agree, that opinion becomes a part of our belief system.

Little by little, all these opinions modify our behavior, and in our mind we form an image of ourselves according to what other people say we are: ``I'm beautiful; I'm not so beautiful.

I'm smart; I'm not so smart.

I'm a winner; I'm a loser.

I'm good at this; I'm bad at that.''

%
At a certain point, all the opinions of our parents and teachers, religion and society, make us believe that we need to be a certain way in order to be accepted.

They tell us the way we \textit{should} be, the way we \textit{should} look, the way we \textit{should} behave.

We need to be \textit{this} way; we shouldn't be \textit{that} way - and because it's not okay for us to be what we are, we start pretending to be what we are not.

The fear of being rejected becomes the fear of not being good enough, and we start searching for something that we call \textit{perfection}.

In our search, we form an image of perfection, the way we wish to be, but we know that we are not, and we begin to judge ourselves for that.

We don't like ourselves, and we tell ourselves, ``Look how silly you look, how ugly you are.

Look how fat, how short, how weak, how stupid you are.''

This is when the drama begins, because now the symbols are going against us.

We don't even notice that we've learned to use the symbols to reject ourselves.

%
Before domestication, we don't care what we are or what we look like.

\textit{Our tendency is to explore, to express our creativity, to seek pleasure and avoid pain}.

\textit{As little children, we are wild and free}; we run around naked without self-consciousness and selfjudgment.

We speak the truth because we live in truth.

\textit{Our attention is in the moment; we are not afraid of the future or ashamed of the past}.

After domestication, we try to be good enough for everybody else, but we are no longer good enough for ourselves, because we can never live up to our image of perfection.

%
All of our normal humans tendencies are lost in the process of domestication, and we begin to search for what we have lost.

We start searching for freedom because we are no longer free to be what we really are; we start searching for happiness because we are no longer happy; we start searching for beauty because we no longer believe that we are beautiful.

%
We continue to grow, and in our adolescence, our body is programmed to introduce a substance we call \textit{hormones}.

Our physical body is no longer a child's, and we don't fit in with the way of life we lived before.

We don't want to hear our parents tell us what to do and what not to do.

We want our freedom; we want to be ourselves, but we are also afraid to be by ourselves.

People tell us, ``You're not a child anymore,'' but we're not an adult either, and it's a difficult time for most humans.

By the time we are teenagers, we don't need anyone to domesticate us; we have learned to judge ourselves, punish ourselves, and reward ourselves according to the same belief system we were given, and using the same system of punishment and reward.

The domestication may be easier for people in some places in the world, and harder for people in other places, but \textit{in general none of us has a chance of escaping the domestication}.

None os us.

Finally, the body matures and everything changes again.

We start searching once again, but now, more and more, what we are searching for is our \textit{self}.

We are searching for love because we have learned to believe that love is somewhere outside of us; we are searching for justice because there is no justice in the belief system we were taught; we are searching for truth because we only believe in the knowledge we have stored in our mind.

And of course, we're still searching for perfection because now we agree with the rest of the humans that ``nobody's perfect.''

\section{Symbols \& Agreements: The Art of Humans}
\textsc{During all the years that we grow up, we make} countless agreements with ourselves, with society, with everybody around us.

But the most important agreements are the ones we make with ourselves by understanding the symbols we learned.

The symbols are telling us what we believe about ourselves; they're telling us what we are and what we are not, what is possible and what is not possible.

\textit{The voice of knowledge is telling us everything that we know, but who tells us if what we know is the truth?}

%
When we go to grammar school, high school, and college, we acquire a lot of knowledge, \textit{but what do we really know?}

\textit{Do we master the truth?}

No, we master a language, a symbology, and that symbology is only the truth because we \textit{agree}, not because it's \textit{really} the truth.

Wherever we are born, whatever language we learn to speak, we find that almost everything we know is really about agreements, beginning with the symbols that we learn.

%
If we are born in England, we learn English symbols.

If we are born in China, we learn Chinese symbols.

But whether we learn English, Chinese, Spanish, German, Russian, or any other language, the symbols only have value because we assign them a value and agree on their meaning.

If we don't agree, the symbols are meaningless.

The word \textit{tree}, e.g., is meaningful for people who speak English, but ``tree'' doesn't mean anything unless we \textit{believe} that it means something, unless we \textit{agree}.

What it means for you, it means for me, and that's why we understand one another.

What I'm saying right now you understand because we agree with the meaning of every word that was programmed in our mind.

But this doesn't mean that we completely agree.

Each of us gives a meaning to every word, and it's not exactly the same for everyone.

%
If we focus our attention on the way any word is created, we find that whatever meaning we give to a word is there for no real reason.

We put words together from nowhere; we make them up.

Humans invent every sound, every letter, every graphic symbol.

We hear a sound like ``A'' and we say, ``This is the symbol for that sound.''

We draw a symbol to represent the sound, we put th symbol and the sound together, and we give it a meaning.

Then every word in our mind has a meaning, but not because it's real, not because it's truth.

It's just an agreement with ourselves, and with everybody else who learns the same symbology.

%
If we travel to a country where people speak a different language, we suddenly realize the importance and power of agreement.

Un \'arbol es s\'olo un \'arbol, el sol es s\'olo el sol, la tierra es s\'olo la tierra si estamos de acuerdo.

\textsf{A sentence in Greek.}

Ein Baum ist nur ein Baum, die Sonne ist nur die Sonne, die Erde ist nur die Erde wenn wir uns darauf verst\"andigt haben.

\textsf{A sentence in Chinese.}

A tree is only a tree, the sun is only the sun, the earth is only the earth if we agree.

These symbols have no meaning in France, Russia, Turkey, Sweden, or in any other place where the agreements are different.

%
If we learn to speak English and we go to China, we hear people talking, but we don't understand a word they are saying.

Nothing makes sense to us because it's not the symbology that we learned.

Many things are foreign to us; it's like being in another world.

If we visit their places of worship, we find that their beliefs are completely different, their rituals are completely different, their mythologies have nothing to do with what we learned.

1 way to understand their culture is to learn the symbols they use, which means their language, but if we learn a new way of being, a new religion or philosophy, this may create a conflict with what we learned before.

New beliefs clash with old beliefs, and the doubt comes right away: ``What is right and what is wrong?

Is it true what I learned before?

Is it true what I'm learning right now?

What is the truth?''

%
The truth is that all of our knowledge, 100\% of it, is nothing more than symbolism or words that we invent for the need to understand and express what we perceive.

Every word in our mind and on this page is just a symbol, but every word has the power of our faith because we \textit{believe} in its meaning without a doubt.

Humans construct an entire belief system made up of symbols; we build an entire edifice of knowledge.

Then we use everything we know, which is nothing but symbology, to justify what we believe, to try to explain 1st to ourselves, then to everybody around us, the way we perceive ourselves, the way we perceive the entire universe.

If we are aware of this, then it's easy to understand that all of the different mythologies, religions, and philosophies around the world, all of the different beliefs and ways of thinking, are nothing but agreements with ourselves and with other humans.

They're \textit{our creation, but are they true?}

Everything that exists is true: the earth is true, the stars are true, the entire universe has always been true.

But the symbols that we use to construct what we know are only true because we say so.
\begin{center}
    $\star$
\end{center}

There's a beautiful story in the Bible that illustrates the relationship between God and humans.

In this story, Adam and God are walking together around the word, and God asks Adam what he wants to name everything.

1 by 1, Adam gives a name to everything he perceives.

``Let's call this a \textit{tree}.

Let's call this a \textit{bird}.

Let's call this a \textit{flower}$\ldots$''

And God agrees with Adam.

The story is about the creation of symbols, the creation of an entire language, and it works by agreement.

It's like \textit{2 sides of the same coin}: We can say that \textit{1 side is pure perception, what Adam perceives; the other side is the meaning that Adam gives to whatever he perceives}.

There's the \textit{object of perception}, which is the truth, and there's our interpretation of the truth, which is just a \textit{point of view}.

\fbox{The truth is objective, and we call it \textit{science}.}

\fbox{Our interpretation of the truth is subjective, and we call it \textit{art}.}

Science and art; the truth, and our interpretation of the truth.

The real truth is life's creation, and it's the absolute truth because it's truth for everyone.

Our interpretation of the truth is our creation, and it's a relative truth because it's only truth by agreement.

With this awareness, we can begin to understand the human mind.

%
All humans are programmed to perceive the truth, and we don't need a language to do this.

But in order to \textit{express} the truth, we need to use a language, and that expression is our art.

It's no longer the truth because words are symbols, and symbols can only represent or ``symbolize'' the truth.

E.g., we can see a tree even if we don't know the symbol ``tree.''

Without the symbol, we just see an object.

The object is real, it is truth, and we perceive it.

Once we call it a \textit{tree}, we are using art to express a point of view.

By using more symbols, we can describe the tree - every leaf, every color.

We can say it's a big tree, a small tree, a beautiful tree, an ugly tree, but is it the truth?

No, the tree is still the same tree.

%
Our interpretation of the tree will depend on our emotional reaction to the tree, and our emotional reaction will depend on the symbols that we use to recreate the tree in our mind.

As you can see, our interpretation of the tree is not exactly the truth.

But our interpretation is a \textit{reflection} of the truth, and that reflection is what we call the \textit{human mind}.

The human mind is nothing but a virtual reality.

It isn't real.

What's real is truth.

What's truth is truth for everyone.

But the virtual reality is our personal creation; it's our art, and it's only ``truth'' for each one of us.

All humans are artists, \textit{all} of us.

Every symbol, every word, is a little piece of art.

From my point of view, and thanks to our programming, our greatest masterpiece of art if the use of a language to create an entire virtual reality within our mind.

The virtual reality we create could be a clear reflection of the truth, or it could be completely distorted.

Either way, it's art.

Our creation could be our personal heaven, or it could be our personal hell.

It doesn't matter; it's art.

But what we can do with the awareness of what is truth and what is virtual is endless.

The truth leads to self-mastery, to a life that's very easy; our distortion of the truth often leads to needless conflict and human suffering.

\fbox{\textit{Awareness makes all the difference}.}

%
Humans are born with awareness; we are born to perceive the truth, but we accumulate knowledge, and we learn to deny what we perceive.

We practice not being aware, and we master not being aware.

\textit{The word is pure magic}, and we learn to use our magic against ourselves, against creation, against our own kind.

To be aware means to open our eyes to see the truth.

When we see the truth, we see everything just as it is, not the way we believe it is, not the way we wish it to be.

Awareness opens the door to millions of possibilities, and if we know that we are the artist of our own life, we can make a choice from all those possibilities.

%
What I'm sharing with you comes from my personal training, which I call \textit{Toltec Wisdom}.

\textit{Toltec} is a N\'ahuatl word meaning \textit{artist}.

From my point of view, to be a \textit{Toltec} has nothing to do with any philosophy or place in the world.

To be a Toltec is just to be an artist.

A Toltec is an artist of the spirit, and as artists we like beauty; we don't like what is not beauty.

If we become better artists, our virtual reality becomes a better reflection of the truth, and we can create a masterpiece of heaven with our art.

%
Thousands of years ago, the Toltec created 3 masteries of the artist: \textit{the mastery of awareness, the master of transformation}, and \textit{the mastery of love, intent}, or \textit{faith}.

The separation is just for our understanding, because the 3 masteries become only one.

The truth is only one, and the truth is what we are talking about.

These 3 masteries guide us out of suffering and return us to our true nature, which is happiness, freedom, and love.

%
The Toltec understood that we are going to create a virtual reality with or without awareness.

If it's with awareness, we're going to enjoy our creation.

And whether we facilitate the transformation or resist it, our virtual reality is always transforming.

If we practice the art of transformation, soon we're facilitating the transformation, and instead of using our magic against ourselves, we are using our magic for the expression of our happiness and our love.

When we master love, intent, or faith, we master the dream of our life, and when all 3 masteries are accomplished, we reclaim our divinity and become one with God.

This is the goal of the Toltec.

%
The Toltec didn't have the technology that we have at the present time; they didn't know about the virtual reality of computers, but they knew how to master the virtual reality of the human mind.

\textit{The mastery of the human mind requires complete control of the attention - the way we interpret and react to information we perceive from inside of us and outside of us}.

The Toltec understood that each one of us is just like God, but instead of creating, we re-create.

\textit{And what do we re-create?}

\textit{What we perceive}.

That is what becomes the human mind.

%
\textit{If we can understand what the human mind is, and what the human mind goes, we can begin to separate reality from virtual reality, or pure perception, which is truth, from symbology, which is art}.

\fbox{\textit{Self-mastery is all about awareness, and it begins with self-awareness}.}

1st to be aware of what is real, and then to be aware of what is virtual, which means what we believe about what is real.

With this awareness, we know that we can change what is virtual by changing what we believe.

\textit{What is real we cannot change, and it doesn't matter what we believe}.

\section{The Story of You: The 1st Agreement: Be Impeccable with Your Word}

\section{Every Mind is a World: The 2nd Agreement: Don't Take Anything Personally}

\section{Truth or Fiction: The 3rd Agreement: Don't Make Assumptions}

\section{The Power of Belief: The Symbol of Santa Claus}

\section{Practice Makes the Master: The 4th Agreement: Always Do Your Best}

%------------------------------------------------------------------------------%

\begin{center}
    \huge\textsc{Part II: The Power of Doubt}
\end{center}

\section{The Power of Doubt: The 5th Agreement: Be Skeptical, but Learn to Listen}

\section{The Dream of the 1st Attention: The Victims}

\section{The Dream of the 2nd Attention: The Warriors}

\section{The Dream of the 3rd Attention: The Masters}

\section{Becoming a Seer: A New Point of View}

\section{The 3 Languages: What Kind of Messenger Are You?}

\section{Epilogue: Help Me to Change the World}

\section{About the Authors}

%------------------------------------------------------------------------------%

\printbibliography[heading=bibintoc]
\end{document}